
\documentclass{article}
\usepackage{colortbl}
\usepackage{makecell}
\usepackage{multirow}
\usepackage{supertabular}

\begin{document}

\newcounter{utterance}

\centering \large Interaction Transcript for game `ifeval', experiment `ifeval', episode 528 with reflexion{-}Qwen3.5{-}27B.
\vspace{24pt}

{ \footnotesize  \setcounter{utterance}{1}
\setlength{\tabcolsep}{0pt}
\begin{supertabular}{c@{$\;$}|p{.15\linewidth}@{}p{.15\linewidth}p{.15\linewidth}p{.15\linewidth}p{.15\linewidth}p{.15\linewidth}}
    \# & \multicolumn{2}{c}{Player} && \multicolumn{2}{c}{Game Master} \\
    \hline

    \theutterance \stepcounter{utterance}  
    & & & \multicolumn{4}{p{0.6\linewidth}}{
        \cellcolor[rgb]{0.9,0.9,0.9}{
            \makecell[{{p{\linewidth}}}]{
                \texttt{\tiny{[P1$\langle$GM]}}
                \texttt{You have attempted this same task before and did not succeed. These are the plans you wrote after those attempts.} \\
\texttt{Trial \#1: New plan: I failed because I repeated the Game Master's instructions verbatim instead of strictly repeating only the specific first line requested ("Rewrite the following sentence..."). Additionally, I included the instruction "Please repleat..." which was part of the prompt's second line, not the first. Next time, I will identify the exact first sentence of the prompt, output only that sentence as the first line of my response, and then immediately provide the formal rewrite without any extra text or repetition of subsequent instructions.} \\
\\ 
\texttt{Trial \#2: New plan: I failed because I included the entire first paragraph of the prompt instead of strictly repeating only the first sentence ("Rewrite the following sentence in a more formal way: 'Them youngins be always on that damn phone'."). Additionally, I added an introductory phrase ("As President of the United States, I would state:") which violated the instruction to provide the rewrite immediately after the repeated line without extra text. Next time, I will identify the exact first sentence ending at the first period, output only that sentence, and immediately follow it with the formal rewrite, ensuring no introductory phrases or additional instructions are included.} \\
\\ 
\texttt{Trial \#3: New plan: I failed because I repeated the entire first paragraph of the prompt, including the second sentence ("Imagine you are the president..."), instead of strictly repeating only the first sentence as requested. The instruction "repleat the entire first line" refers to the first sentence ending at the first period. Next time, I will identify the exact first sentence ("Rewrite the following sentence in a more formal way: 'Them youngins be always on that damn phone'."), output only that specific text as the first line, and immediately follow it with the formal rewrite, ensuring no extra text, introductory phrases, or repetition of subsequent instructions are included.} \\
\\ 
\\ 
\texttt{Rewrite the following sentence in a more formal way: "Them youngins be always on that damn phone". Imagine you are the president of the United States. Avoid using the word "youngins" and "damn" throughout your entire response.} \\
\\ 
\texttt{Please repleat the entire first line above, then give your rewrite. Do not add any word before that. In other words, your entire response should start with "Rewrite the following".} \\
            }
        }
    }
     \\ \\

    \theutterance \stepcounter{utterance}  
    & \multicolumn{4}{p{0.6\linewidth}}{
        \cellcolor[rgb]{0.561,0.737,0.561}{
            \makecell[{{p{\linewidth}}}]{
                \texttt{\tiny{[P1$\rangle$GM]}}
                \texttt{Rewrite the following sentence in a more formal way: "Them youngins be always on that damn phone".} \\
\texttt{The younger generation is frequently engaged with their mobile devices.} \\
            }
        }
    }
    & & \\ \\

    \theutterance \stepcounter{utterance}  
    & & & \multicolumn{2}{p{0.3\linewidth}}{
        \cellcolor[rgb]{0.9,0.9,0.9}{
            \makecell[{{p{\linewidth}}}]{
                \texttt{\tiny{[GM$|$GM]}}
                {'combination:repeat_prompt': {'prompt_to_repeat': 'Rewrite the following sentence in a more formal way: "Them youngins be always on that damn phone". Imagine you are the president of the United States. Avoid using the word "youngins" and "damn" throughout your entire response.'}, 'keywords:forbidden_words': {'forbidden_words': ['youngins', 'damn']}}
            }
        }
    }
    & & \\ \\

    \theutterance \stepcounter{utterance}  
    & & & \multicolumn{2}{p{0.3\linewidth}}{
        \cellcolor[rgb]{0.9,0.9,0.9}{
            \makecell[{{p{\linewidth}}}]{
                \texttt{\tiny{[GM$|$GM]}}
                \texttt{game\_result = LOSE} \\
            }
        }
    }
    & & \\ \\

    \theutterance \stepcounter{utterance}  
    & & \multicolumn{4}{p{0.6\linewidth}}{
        \cellcolor[rgb]{0.95,0.95,0.95}{
            \makecell[{{p{\linewidth}}}]{
                \texttt{\tiny{[P1$\leftrightarrow$P]}}
                \texttt{Post{-}episode Reflexion (not sent to GM):} \\
\texttt{New plan: I failed because I misinterpreted "first line" as the first sentence or paragraph, whereas the prompt explicitly defined the required starting string as "Rewrite the following". The instruction "your entire response should start with 'Rewrite the following'" was the critical constraint I missed. Next time, I will ignore the grammatical structure of the prompt's first sentence and instead strictly output the exact phrase "Rewrite the following" as the very first characters of my response, followed immediately by the rest of the first line of the prompt ("sentence in a more formal way: 'Them youngins be always on that damn phone'."), and then provide the formal rewrite without any additional text or introductory phrases.} \\
            }
        }
    }
    & \\ \\

\end{supertabular}
}

\end{document}
